\section{Conclusion}\label{sec:conclusion}

Le système réalisé couvre l'ensemble du cahier des charges : la base de données modélise toutes les fonctionnalités demandées et en garantit les règles métier par des contraintes, des déclencheurs et des procédures stockées, et l'application web dépasse le minimum exigé en exposant, en plus de l'authentification, des stations, des locations et des statistiques, la gestion des abonnements, des signalements, de la flotte et des comptes.

Le passage à un jeu de données réaliste de plusieurs gigaoctets a été l'étape la plus formatrice : il a révélé des requêtes correctes mais inexploitables à ce volume, et imposé une démarche rigoureuse de mesure, de lecture des plans d'exécution et de réécriture, en vérifiant à chaque fois que les résultats restaient identiques.

\afaire{bilan personnel de l'équipe, limites (statistiques sur de longues périodes, maintenance et régulation partiellement exposées dans l'interface) et perspectives (tables d'agrégats, carte des stations, application mobile\ldots).}
